\scanpage{70}
\mspage{9}
métriques euclidiennes \uline{canoniques} sur
les $T_{X,a_i}$, décrivant leur structure
conforme. Dans le cas $g\geq 2$ et $=$
encore pour $g=1$, cela ne semble pas
inquiétant, puisqu'on a une
métrique canonique sur $X$. Par
contre dans le cas $g=0$, cela semble
étrange, puisque dans ce cas on n'a
\uline{pas} de métrique canonique sur $X$.

\uline{12.1.}

5. \uline{Retour sur le principe de chirurgie conforme}

La formulation des lemmes \inserted{locaux} 2 et 3, qui
sont l'ingrédient technique essentiel pour le
th. 4 (et lui sont équivalents), est un peu
alambiquée. Il est plus élégant de se
réduire à des lemmes globaux de type
uniformisation par le disque $\mathbb{D}$. Cela
donne le
% archive sheet number: 72
